\subsection{\textbf{Sprint 1}}

Na Sprint 1, o time foi dividido em \textit{squads} de quatro pessoas. No entanto, essa divisão não foi feita de maneira estratégica, pois as \textit{tasks} já haviam sido atribuídas previamente. Com isso, as \textit{squads} acabaram tendo um papel meramente protocolar. Como eu já havia me aprofundado bastante no estudo da biblioteca Three.js, assumi a \textit{task} de criar o mapa para a tela do visitante. Estava confiante de que conseguiria entregar um resultado sólido, e mesmo que fosse um mapa estático e não dinâmico, já representaria um bom avanço para esta etapa do projeto.

Durante as aulas, envolvi-me bastante com a equipe que pesquisava formas de automatizar a geração de mapas a partir de arquivos DWG ou DXF. Infelizmente, até a data da apresentação para o \textit{stakeholder}, não conseguimos evoluir significativamente nessa frente.

Com uma semana de antecedência em relação à apresentação, comunicamos a todos que o \textit{code freeze} seria realizado dois dias antes da aula, a fim de garantir tempo hábil para revisar cuidadosamente cada \textit{pull request} e resolver possíveis conflitos de forma organizada.

Concluí minha \textit{task} com certa antecedência e, ao verificar o andamento do \textit{backend}, percebi que nenhum outro \ac{ages} III havia contribuído com o repositório até então. O código encontrava-se desorganizado e sem uma estrutura adequada, mas, por estarmos muito avançados na \textit{sprint}, não havia mais tempo hábil para reestruturar a arquitetura.

Na véspera do \textit{code freeze}, auxiliei o \ac{ages} II Kauã Rodrigues Souza com dúvidas relativas à sua \textit{task} e aos critérios de entrega esperados. Percebemos que havia ruído na comunicação entre os \ac{ages} III e que não havíamos alinhado previamente como seriam feitos os \textit{pull requests}. Diante do risco que isso representava para a entrega, tomei a iniciativa de revisar e aprovar todos os PRs do repositório do \textit{backend} sozinho. Apesar da responsabilidade, o processo ocorreu de forma tranquila, e fiquei satisfeito com a qualidade do código entregue por aqueles que contribuíram.

Entregamos ao \textit{stakeholder} praticamente todos os itens planejados para a \textit{sprint}, com exceção do componente de barra de pesquisa no \textit{frontend}. Durante a retrospectiva, discutimos pontos de melhoria e definimos ações concretas para evitar os problemas enfrentados nesta etapa. De forma geral, fiquei satisfeito com a entrega, embora tenha ficado com certa desconfiança quanto ao nível de comprometimento de alguns integrantes do time.
